\subsubsection{Specialization, loss masks and output semantics}
\label{sec:output-semantics}
A global output index, a supervised loss term, a specialist output and an evaluated endpoint are different objects. The existence of a logit for a named diagnosis establishes that the implementation can emit that score; it does not establish that this output received direct training support in every stage. Recovered V3F loss indexing supports the stated P8 and P4 objectives, while earlier or different development families cannot be inferred from the released architecture alone. In particular, the separate historical 41M experiments do not document direct supervision for every output of the released 58.35M backbones.

Primary, secondary-candidate and auxiliary are reporting roles in Table~\ref{tab:contracts}. They do not identify mutually exclusive patient groups. A conduction statement such as RBBB may coexist with a rate or rhythm statement, and a pathology reference can coexist with any rhythm. A categorical top-1 view compresses that set into one reported class, whereas thresholded membership can retain several classes or none. P6 and P8 therefore require both an accepted-reference-set definition and a separate membership policy. Calling NSR an arrhythmia or treating it as a residual abnormal class would change the task. A residual Other endpoint must instead define which out-of-profile statements count as positive and which records have genuinely known negative references.

The pathology distinction is equally important. ALMI and ILMI identify different infarct-territory statements from ASMI and IMI; their output availability does not make them evaluated P4 targets. LAE and LAO may be combined only in an explicitly bounded experimental reference, not relabeled as interchangeable clinical diagnoses. Similarly, LQTS is a syndrome label and cannot be validated by automatically substituting a prolonged-QT observation. The later residual experiments are described with their actual target unions, eligible denominators and fitted heads. They do not create a newly validated released Other class.

\subsubsection{Signal scaling and morphology limitations}
\label{sec:signal-limits}
Per-lead standardization makes the numeric input less dependent on recording gain, but removes absolute amplitude and changes cross-lead amplitude relationships. The morphology power ratios and correlation proxies operate on that prepared representation. They cannot recover the discarded physical millivolt scale. This matters for LVH and other voltage-dependent findings, and motivates testing a future model that retains verified physical amplitudes alongside normalized waveforms. The present results do not isolate normalization as the cause of LVH error, because no controlled amplitude-preserving ablation was performed.

The peak detector and pre-peak lag estimator are heuristic calculations. A maximum deflection before a detected peak need not be a P wave, and the thresholded width need not correspond to expert QRS onset and offset. Filtering, resampling, sampling phase, lead choice and artifact can all change the derived vector. A lead-validity mask describes the input contract but does not certify adequate signal quality. The known flat-signal failure is therefore reported as a methodological limitation rather than silently repaired before replay. Prospective use would require explicit quality criteria, interval-reference validation and failure handling evaluated on representative recordings.

Classical QRS detection provides a signal-processing context for automated peak localization \citep{Pan1985Real}. The present extractor is defined by its own recovered thresholds and timing rules in Table~\ref{tab:morphology}. No timing-annotation study was performed to establish its peak or interval accuracy. Peak detection, interval measurement and ECG-statement classification must therefore remain separate claims.

The reconstructed runtime audit separates deterministic same-input repetition from reproduction of the June calculation. Repeating a forward pass on the same prepared batch tests one numerical configuration. It does not reconstruct unavailable historical waveform windows, library versions, kernel choices, import paths or masks. Sensitivity to preprocessing and device arithmetic must therefore be considered alongside threshold agreement. Stable threshold decisions on one matched panel can coexist with non-equivalent scores and with changed decisions on other ECGs.
